\chapter*{Abstract}
\addcontentsline{toc}{chapter}{Abstract}

This thesis investigates the use of Large Language Models (LLMs) for high-level planning and
adaptive replanning in heterogeneous multi-drone systems. The system uses a hierarchical
architecture in which a cloud-based planning pipeline generates mission plans from natural-language
instructions, while a smaller locally deployable model performs task admission based on the 
drone's state.

The planning pipeline decomposes a human-language mission description into atomic subtasks, assigns them to drones
according to their capabilities, and generates an execution schedule. Before and during execution, the task 
admission module determines whether assigned subtasks are feasible given the drone's state. 
If a drone failure or task rejection occurs, the affected subtasks are returned to the cloud-based planner for rescheduling.

The system was evaluated in three parts. First, several cloud-based LLMs were compared for task
decomposition, allocation, and scheduling using missions with one to ten subtasks. The generated
schedules were evaluated against a vehicle routing problem baseline in terms of schedule quality,
validity, and inference time. GPT-5-mini was selected for the cloud-based planner because it provided
the best trade-off between quality, latency, and cost. Second, seven quantized locally
deployable models were evaluated for task admission using scenarios involving drone-state errors, 
link-quality errors, and flight-time errors. The qwen3:1.7b model achieved
the highest reliability and was selected for the local task admission module.

Finally, the selected components were integrated into a simulation environment containing
heterogeneous drones and diverse target objects, while also simulating battery drainage,
link-quality variation, and drone failures. The simulation case study demonstrated that the system
can generate mission schedules, perform task admission checks, detect execution-time failures,
reschedule affected subtasks, and continue execution with the remaining drones. The physical Anafi
drone platform and laboratory integration setup are also described to relate the symbolic planner
skills to real drone capabilities.

The results show that LLMs can support flexible high-level mission planning and replanning when
combined with structured task representations, capability-based allocation, and execution-time
feasibility checks. However, the complete closed-loop architecture was evaluated in simulation rather
than deployed on physical drones, leaving full real-world validation as future work.